\begin{figure}
    \centering
    \includegraphics[width=0.95\textwidth]{figures/overview}
    \caption{
        System Overview. (a) We use automatic domain randomization (ADR) to generate a growing distribution of simulations with randomized parameters and appearances. We use this data for both the control policy and vision-based state estimator. (b) The control policy receives observed robot states and rewards from the randomized simulations and learns to solve them using a recurrent neural network and reinforcement learning. (c) The vision-based state estimator uses rendered scenes collected from the randomized simulations and learns to predict the pose as well as face angles of the Rubik's cube using a convolutional neural network (CNN), trained separately from the control policy. (d) To transfer to the real world, we predict the Rubik's cube's pose from 3 real camera feeds with the CNN and measure the robot fingertip locations using a 3D motion capture system. The face angles that describe the internal rotational state of the Rubik's cube are provided by either the same vision state estimator \emph{or} the Giiker cube,  a custom cube with embedded sensors and feed it into the policy network.
    }
    \label{fig:overview}
\end{figure}

Building robots that are as versatile as humans remains a grand challenge of robotics. While humanoid robotics systems exist~\cite{bostondynamicsatlas,shadow-robot,toyotathr3,sakagami2002intelligent,asfour2013armar}, using them in the real world for complex tasks remains a daunting challenge. Machine learning has the potential to change this by \emph{learning} how to use sensor information to control the robot system appropriately instead of hand-programming the robot using expert knowledge.

However, learning requires vast amount of training data, which is hard and expensive to acquire on a physical system.  Collecting all data in simulation is therefore appealing. However, the simulation does not capture the environment or the robot accurately in every detail and therefore we also need to solve the resulting sim2real transfer problem.
Domain randomization techniques~\cite{tobin2017domain, peng2017sim} have shown great potential and have demonstrated that models trained only in simulation can transfer to the real robot system.

In prior work, we have demonstrated that we can perform complex in-hand manipulation of a block~\cite{openai2018learning}. This time, we aim to solve the manipulation and state estimation problems required to solve a Rubik's cube with the Shadow Dexterous Hand~\cite{shadow-robot} using only simulated data. This problem is much more difficult since it requires significantly more dexterity and precision for manipulating the Rubik's cube. The state estimation problem is also much harder as we need to know with high accuracy what the pose and internal state of the Rubik's cube are. We achieve this by introducing a novel method for automatically generating a distribution over randomized environments for training reinforcement learning policies and vision state estimators. We call this algorithm \emph{automatic domain randomization} (ADR). We also built a robot platform for solving a Rubik's cube in the real world in a way that complements our machine learning approach. \autoref{fig:overview} shows an overview of our system.

We investigate why policies trained with ADR transfer so well from simulation to the real robot. We find clear signs of emergent learning that happens at \emph{test time} within the recurrent internal state of our policy. We believe that this is a direct result of us training on an ever-growing distribution over randomized environments with a memory-augmented policy. In other words, training an LSTM over an ADR distribution is implicit meta-learning. We also systematically study and quantify this observation in our work.


The remainder of this manuscript is structured as follows. \autoref{sec:task} introduces two manipulation tasks we consider here. \autoref{sec:physical-setup} describes our physical setup and \autoref{sec:sim} describes how our setup is modeled in simulation. We introduce a new algorithm called automatic domain randomization (ADR), in \autoref{sec:adr}. In \autoref{sec:policy} and \autoref{sec:vision} we describe how we train control policies and vision state estimators, respectively. We present our key quantitative and qualitative results on the two tasks in \autoref{sec:exp-adr}. In \autoref{sec:exp-meta} we systematically analyze our policy for signs of emergent meta-learning. \autoref{sec:related} reviews related work and we conclude with \autoref{sec:conclusion}.

If you are mostly interested in the machine learning aspects of this manuscript, \autoref{sec:adr},  \autoref{sec:policy},  \autoref{sec:vision}, \autoref{sec:exp-adr}, and \autoref{sec:exp-meta} are especially relevant. If you are interested in the robotics aspects, \autoref{sec:physical-setup}, \autoref{sec:sim}, and \autoref{sec:result-rubik} are especially relevant.
